% arara: xelatex
\documentclass[12pt]{article}

% \usepackage{physics}

\usepackage{hyperref}
\hypersetup{
    colorlinks=true,
    linkcolor=blue,
    filecolor=magenta,      
    urlcolor=cyan,
    pdftitle={Overleaf Example},
    pdfpagemode=FullScreen,
    }

\usepackage{tikzducks}

\usepackage{tikz} % картинки в tikz
\usetikzlibrary{shapes, arrows, positioning}
\usepackage{microtype} % свешивание пунктуации

\usepackage{array} % для столбцов фиксированной ширины

\usepackage{indentfirst} % отступ в первом параграфе

\usepackage{sectsty} % для центрирования названий частей
\allsectionsfont{\centering}

\usepackage{amsmath, amsfonts, amssymb} % куча стандартных математических плюшек

\usepackage{comment}

\usepackage[top=2cm, left=1.2cm, right=1.2cm, bottom=2cm]{geometry} % размер текста на странице

\usepackage{lastpage} % чтобы узнать номер последней страницы

\usepackage{enumitem} % дополнительные плюшки для списков
%  например \begin{enumerate}[resume] позволяет продолжить нумерацию в новом списке
\usepackage{caption}

\usepackage{url} % to use \url{link to web}

\usepackage{xcolor}   % для цвета текста
\usepackage{ulem}     % для зачёркивания (\sout)

\usepackage[symbol]{footmisc}

\newcommand{\smallduck}{\begin{tikzpicture}[scale=0.3]
    \duck[
        cape=black,
        hat=black,
        mask=black
    ]
    \end{tikzpicture}}

\usepackage{fancyhdr} % весёлые колонтитулы
\pagestyle{fancy}
\lhead{Stochastic Processes}
\chead{}
\rhead{Final Exam, 2025-12-26}
\lfoot{In this exam $(W_t)$ always denotes a standard Wiener process.}
\cfoot{}
\rfoot{}

\renewcommand{\headrulewidth}{0.4pt}
\renewcommand{\footrulewidth}{0.4pt}

\usepackage{tcolorbox} % рамочки!

\usepackage{todonotes} % для вставки в документ заметок о том, что осталось сделать
% \todo{Здесь надо коэффициенты исправить}
% \missingfigure{Здесь будет Последний день Помпеи}
% \listoftodos - печатает все поставленные \todo'шки


% более красивые таблицы
\usepackage{booktabs}
% заповеди из докупентации:
% 1. Не используйте вертикальные линни
% 2. Не используйте двойные линии
% 3. Единицы измерения - в шапку таблицы
% 4. Не сокращайте .1 вместо 0.1
% 5. Повторяющееся значение повторяйте, а не говорите "то же"


\setcounter{MaxMatrixCols}{20}
% by crazy default pmatrix supports only 10 cols :)


\usepackage{fontspec}
\usepackage{libertine}
\usepackage{polyglossia}

\setmainlanguage{russian}
\setotherlanguages{english}


\usetikzlibrary{automata, positioning, arrows.meta, calc, shapes.multipart}


% download "Linux Libertine" fonts:
% http://www.linuxlibertine.org/index.php?id=91&L=1
% \setmainfont{Linux Libertine O} % or Helvetica, Arial, Cambria
% why do we need \newfontfamily:
% http://tex.stackexchange.com/questions/91507/
% \newfontfamily{\cyrillicfonttt}{Linux Libertine O}

\AddEnumerateCounter{\asbuk}{\russian@alph}{щ} % для списков с русскими буквами
% \setlist[enumerate, 2]{label=\asbuk*),ref=\asbuk*}

%% эконометрические сокращения
\DeclareMathOperator{\Cov}{\mathbb{C}ov}
\DeclareMathOperator{\Corr}{\mathbb{C}orr}
\DeclareMathOperator{\Var}{\mathbb{V}ar}
\DeclareMathOperator{\col}{col}
\DeclareMathOperator{\row}{row}

\let\P\relax
\DeclareMathOperator{\P}{\mathbb{P}}

\DeclareMathOperator{\E}{\mathbb{E}}
% \DeclareMathOperator{\tr}{trace}
\DeclareMathOperator{\card}{card}

\DeclareMathOperator{\Convex}{Convex}
\DeclareMathOperator{\plim}{plim}

\newcommand{\cN}{\mathcal{N}}
\newcommand{\cF}{\mathcal{F}}

\newcommand{\RR}{\mathbb{R}}
\newcommand{\NN}{\mathbb{N}}
\newcommand{\hb}{\hat{\beta}}
\newcommand{\dPois}{\mathrm{Pois}}
\newcommand{\dUnif}{\mathrm{Unif}}





\begin{document}

Сегодня по народному календарю  Евстратиев день. 
Нельзя в этот день сквернословить — не то ведьмы с неба упадут прямо на голову. 
Также нельзя веник на крыльце оставлять — не то ведьмы утащат.

\begin{enumerate}
\item The Legendary November Midterm Grader grades tests according to Poisson point process with variable intensity rate.
The first 5 weeks the rate is equal to $\lambda_1 = 1$ test per week. 
Later the rate increases to $\lambda_2 = 100$ tests per week. 
Let $X_t$ be the number of tests graded after $t$ weeks. 

\begin{enumerate}
\item {[2]} Sketch the curve $v(t) = \Var(X_t)$.
\item {[4]} Using normal approximation find the probability that LNMG will check not less than $120$ works at time $t=6$.
\item {[4]} Now $t=4.5$ weeks have already passed. Find the expected time for the next test checked.
\end{enumerate}
Note: in point (b) you may use the standard normal cdf $F()$ in your answer.

%Packs of assignments arrive independently as $P_t \sim \dPois(\lambda_1)$ with intensity $\lambda_1 = 2$ packs per week. 
%Grading a single pack is a creative and randomly timed process. 
%The grading rate is also independent of arrivals and number of graded packs is $G_t \sim \dPois(\mu_1)$ with intensity $\mu_1 = 1$ pack per 3 weeks. 
%At the beginning there are 3 ungraded packs. 
%Let $X_t$ be the total number of ungraded packs at time $t \geq 0$ (including the initial packs).

%\begin{enumerate}
%    \item {[5]} Find $\E(X_t)$ and $\Var(X_t)$. 
%    What happens to $\E(X_t)$ as $t \to \infty$? 
%    Provide a practical interpretation of this result in the context of the Professor's workload.
    
%    \item {[5]} Suppose Professor Boris realizes his responsibility and suddenly increases his grading rate to $\mu_2 = 2$ packs per week. 
%    No new assignments arrive after this point ($\lambda_2 = 0$). 
%    Let $X_{13}=1$ be the number of ungraded packs. Let $T$ be the time needed to grade all the remaining packs. 
%    Derive the conditional distribution of $T$ given that $X_{13}=1$ and find $\E(T \mid X_{13}=1)$.
% \end{enumerate}

\item The joint distribution of $X$ and $Y$ is given by the table:
\begin{center}
\begin{tabular}{c|ccc}
 & $Y = -1$ & $Y = 0$ & $Y = 1$ \\
\hline
$X = -1$ & $1/8$ & $1/16$ & $1/8$ \\
$X = 0$ & $1/16$ & $1/4$ & $1/16$ \\
$X = 1$ & $1/8$ & $1/16$ & $1/8$ \\
\end{tabular}
\end{center}
Three students are building statistical models with different information sets:
\begin{itemize}
\item Roman Bokhyan: $\cF_1 = \sigma(X, Y)$, the $\sigma$-algebra generated by $X$ and $Y$. 
\item Mikhail N: $\cF_2 = \sigma(X + Y, X - Y)$.
\item Andrey L: $\cF_3 = \sigma(X^2 + Y^2)$.
\end{itemize}
\begin{enumerate}
    \item {[3]} For each pair of $\sigma$-algebras above determine whether one is contained in the other, they are equal, or neither. 
    Justify your claims.
    
    \item {[2]} Andrey and Mikhail decide to share their information. 
    Compute the number of elements in the smallest $\sigma$-algebra containing both $\cF_2$ and $\cF_3$. 
    
    \item {[5]} Help Andrey and compute $\E(X \mid \cF_3)$, $\Var(X \mid \cF_3)$.    
\end{enumerate}

\item The process $S_n$ models the «accumulated wrath of the heavens». Let $S_0 = 0$. 
Each hour independently of previous history $S_n$ increases by $1$ with probability $1/3$, decrease by $1$ with probability $1/6$ or does not change.

The witches fall from the sky when $S_n$ first reaches the level of $4$. 
Let $T = \min\{n:S_n = 4\}$.

\begin{enumerate}
    \item {[4]} For which $\alpha \neq 1$ is the process $M_n = \alpha^{S_n}$ a martingale?
    
    \item {[2]} Check whether $Q_n = S_n - n / 6$ is a martingale.
    
    \item {[4]} Find $\E(T)$.
\end{enumerate}
Hint: you may use Doob's theorem without checking technical details. 

\newpage
\item Consider three processes, 
\[
A_t = \int_0^t W_u^2 \, dW_u, \quad B_t = \int_0^t W_u^4 \, dW_u, \quad C_t = \int_0^t W_u^2 \, du.
\]
\begin{enumerate}
    \item {[2 + 2]} Find $\E(A_t)$, $\E(C_t)$.
    \item {[2 + 2 + 2]} Find $\Var(A_t)$, $\Cov(A_t, B_t)$, $\Var(C_t)$.
\end{enumerate}

\item Be brave!
\begin{enumerate}
\item {[2 + 3]} Find $\E(W_3^4)$, $\E(W_3^4 \mid W_2 = 10)$.
\item {[5]} Find $\Cov(W_1 W_2, W_5 W_6)$.
\end{enumerate}

\item Consider the process
\[
R_t = \exp\left(\int_0^t u^3 dW_u + t^2\right).
\]
\begin{enumerate}
\item {[6]} Write $R_t$ as a sum of a constant, an Ito's integral and a Riemann integral.
\item {[2]} Is $(R_t)$ a martingale?
\item {[2]} Provide an ordinary (non-stochastic) differential equation for $h(t) = \E(R_t)$. 
You don't need to solve this equation. 
\end{enumerate}
Hint: you may write $R_t$ as $R_t = \exp(Q_t)$ and apply Ito's lemma :)


\end{enumerate}


\end{document}

